\scanpage{20}
\textcircled{B}
Supposons à la fois
\[
\begin{aligned}
\alpha)&\quad \widetilde H_0\ \text{est connexe et simplement connexe}\\
\beta)&\quad \pi\subset\widetilde H_0\\
\gamma)&\quad \pi,N,\unsure{\mathfrak{S}|\mathcal{G}|\mathfrak{G}}\ \text{discret}
\end{aligned}
\]
Comme \(\mathfrak{S}\) discret, \(\widetilde H_0\) connexe, on a
\[
\boxed{\left\{\begin{aligned}
\widetilde H_0&=\widetilde H^{0}\quad\text{(composante neutre)}\\
\mathfrak{S}&=\pi_0\widetilde H
\end{aligned}\right.}
\]
Comme \(\pi\subset\widetilde H_0\), on a que \(H=\widetilde H/\pi\) est une extension de \(\mathfrak{S}\) par \(H^0=\widetilde H_0/\pi\)
\[
\boxed{H^0=\widetilde H_0/\pi,\quad \mathfrak{S}=\pi_0(H)}
\]
Comme \(\widetilde H_0\) est simplement connexe, \deleted{c'est le}\inserted{\(\pi\) discret, \(\widetilde H_0\) est} le revêtement universel de \(H^0=\widetilde H_0/\pi\)
\[
\boxed{\left\{\begin{aligned}
\widetilde H^{0}&=\text{rev. universel de }H^0\\
\pi&=\pi_1(H^0)=\pi_1(H)
\end{aligned}\right.}
\]
Donc la situation équivaut :

\(1^\circ)\) donnée d'un groupe top \(H\), \deleted{\unsure{entre}}\unsure{q|i.e.} [qui est donc ext. d'un grpe discret \(\mathfrak{S}=\pi_0(H)\) par \deleted{\unsure{l'un}} groupe connexe \(H^0\)] ; ce dernier \deleted{\unsure{\(2^\circ)\)}} a un groupe \unsure{revêtmt|revêtant} universel \(\widetilde H^0\), extension de \(H^0\) par \(\pi\simeq\pi_1(H^0)\)]
\margin{groupe \uuline{comm.} sur lequel \(\mathfrak{S}\) opère}

\(2^\circ)\) Une extension \(\widetilde H\) de \(H\) par \(\pi\) qui prolonge l'extension \(\widetilde H^0\) de \(H^0\) par \(\pi=\pi_1(H^0)\)

\uline{NB} Une telle extension n'est pas nécessairement unique à isom près ; elle est unique modulo une extension de \(\mathfrak{S}\) par \(\pi\), \inserted{i.e. mod. \unsure{éltmt|élémt|action} de \(H^2(\mathfrak{S},\pi)\),} et sauf erreur l'obstruction à l'existence d'une telle ext. est dans \(H^3(\mathfrak{S},\pi)\). Une telle extension existe si \(H\) est produit semi direct \(\mathfrak{S}.H^0\), car on peut prendre alors \(\widetilde H=\mathfrak{S}.\widetilde H^0\).

% slanted
\margin{\uline{NB} \(G_0=\pi_0\widetilde G=\widetilde G/\widetilde G^0\) (\(\widetilde G^0\to\widetilde H^0\) est un isom) \deleted{\unsure{en termes de}} et comme \(G=\widetilde G/\pi\), on trouve que}
% slanted; same marginal note, branches below its horizontal sequence
\diagram{\begin{tikzcd}[ampersand replacement=\&,column sep=1.2em,row sep=.65em]
\pi\arrow[r]\arrow[dddr] \& \pi_0\widetilde G\arrow[r]\arrow[d,equal] \& \pi_0G\arrow[r] \& 1\\
\&G_0\arrow[d,phantom,"{\cup}"]\&\&\\
\&N\arrow[d,phantom,"{\cup}"]\&\&\\
\&\mathfrak{z}(N)\&\&
\end{tikzcd}}

\(3^\circ)\) Le groupe discret \(N\)

Avec ces données, on trouve que \uline{la catégorie des extensions de \(H\) par \(N\)} équivaut à celle des couples \((\deleted{\unsure{H}}G_0,\varphi)\), où \(\deleted{\unsure{H}}G_0\) est une ext. de \(\mathfrak{S}\) par \(N\) et \(\varphi:\pi\longrightarrow\mathfrak{z}(N)^{\mathfrak{S}}\) un hom \((\mathfrak{z}(N)=\text{centre de }N)\)
% archive sheet number: 19

